\documentclass[11pt,a4paper]{article}
\usepackage{fontspec,amsmath,amssymb,amsthm,microtype}
\setmainfont{lmroman10-regular.otf}[BoldFont=lmroman10-bold.otf,ItalicFont=lmroman10-italic.otf,BoldItalicFont=lmroman10-bolditalic.otf]
\usepackage{xcolor}
\usepackage[most]{tcolorbox}
\usepackage{mathrsfs}
\providecommand{\lt}{<}
\providecommand{\gt}{>}
\setmonofont{DejaVuSansMono.ttf}[BoldFont=DejaVuSansMono-Bold.ttf,ItalicFont=DejaVuSansMono-Oblique.ttf,BoldItalicFont=DejaVuSansMono-BoldOblique.ttf,Scale=0.86]
\newfontfamily\leanfallback{STIXTwoMath-Regular.otf}[Scale=0.86]
\definecolor{leanrule}{gray}{0.62}
\definecolor{leanback}{gray}{0.962}
\definecolor{leantitle}{gray}{0.22}
\definecolor{leanmuted}{gray}{0.38}
\newtcolorbox{leanbox}[1]{enhanced,breakable,colback=leanback,colframe=leanrule,
  boxrule=0.4pt,arc=0.6mm,left=2.4mm,right=2.4mm,top=1.8mm,bottom=1.8mm,
  before skip=0.8em,after skip=1em,fonttitle=\small\bfseries,coltitle=leantitle,
  colbacktitle=leanback,titlerule=0.3pt,title={#1}}
\newlength{\leancharwidth}
\newenvironment{leancode}{\par\smallskip\begingroup\ttfamily\footnotesize\raggedright
  \settowidth{\leancharwidth}{M}\setlength{\parskip}{0pt}\setlength{\parindent}{0pt}
  \frenchspacing}{\par\endgroup\smallskip}
\newcommand{\leanline}[2]{\hangindent=\dimexpr#1\leancharwidth+4\leancharwidth\relax
  \hangafter=1\noindent\hspace*{#1\leancharwidth}#2\par}
\newcommand{\leansource}[1]{\par{\footnotesize\color{leanmuted}#1}\par}
\newcommand{\leanlabel}[1]{\par\smallskip\noindent{\small\bfseries #1}\par}
\usepackage[margin=25mm]{geometry}
\usepackage[unicode,hidelinks]{hyperref}
\hypersetup{pdfauthor={LangevinLean Project}}
\setlength{\emergencystretch}{3em}
\setlength{\parskip}{0.4em}
\setlength{\parindent}{0pt}
\allowdisplaybreaks[2]
\raggedbottom
\title{\texorpdfstring{Reference segments as restrictions of one global solution}{Reference segments as restrictions of one global solution}}
\author{LangevinLean Project}
\date{}
\begin{document}\maketitle
Draft mathematical exposition; these notes may contain mistakes.


We keep the notation of the two preceding sections: \(L^2=L^2(P;\mathbb R^\iota)\), \(\mathcal G[X]=[\nabla U\circ X]\) for an \(L\)-Lipschitz gradient, the Brownian input \(N_a(t)=\sqrt2(\beta(a+t)-\beta(a))\), the chosen segments \(S_n(X)\), the knots \(K_n\) and the reference segments \(Y^{(n)}=S_n(K_n)\) for a step \(h\ge0\) and initial vector \(X_0\in L^2\). Let \(T\in\mathbb R\) and let \(Y:\mathbb R\to L^2\) be a continuous curve satisfying the \emph{global} integral equation


\[Y(t)-X_0=-\int_0^t\mathcal G\bigl(Y(s)\bigr)\,ds+N_0(t),\qquad t\in[0,T].\qquad(\ast)\]


We prove that the segments are pieces of this one curve: \(K_n=Y(nh)\) whenever \(nh\le T\), and \(Y^{(n)}(t)=Y(nh+t)\) for \(t\in[0,h]\) whenever \((n+1)h\le T\). The ingredients are an algebraic identity for the input, \(N_0(a+t)-N_0(a)=N_a(t)\); the observation that the time-shifted curve \(t\mapsto Y(a+t)\) solves the segment equation started at \(Y(a)\) with input \(N_a\); and uniqueness of continuous solutions of the integral equation with given initial value and input. No assertion about laws, stationarity or independence enters. The identification is used later to transfer statements about the law of \(Y(s)\) at fixed times \(s\) to every segment time \(nh+t\).


\section{\texorpdfstring{The increment of the Langevin Brownian input}{The increment of the Langevin Brownian input}}


\textbf{Claim.} Let \(P\) be any measure and \((W^i)_{i\in\iota}\) pre-Brownian motions. For all real numbers \(a\) and \(t\),


\[N_0(a+t)-N_0(a)=N_a(t)\qquad\text{in }L^2.\]


\textbf{Proof.} By definition, \(N_0(s)=\sqrt2(\beta(0+s)-\beta(0))=\sqrt2(\beta(s)-\beta(0))\) for every real \(s\). Hence, using linearity of scalar multiplication,


\[N_0(a+t)-N_0(a)=\sqrt2\bigl(\beta(a+t)-\beta(0)\bigr)-\sqrt2\bigl(\beta(a)-\beta(0)\bigr)=\sqrt2\bigl(\beta(a+t)-\beta(a)\bigr)=N_a(t).\qquad\square\]


No sign condition on \(a\) or \(t\) is needed: the identity is linear algebra in \(L^2\). It is the counterpart for the input of the additivity \(\int_0^{a+t}=\int_0^a+\int_a^{a+t}\) of time integrals, and it gives the input of the equation restarted at time \(a\).

\begin{leanbox}{Lean: \texttt{brownianLangevinNoise\_\allowbreak{}increment}}
\leansource{\texttt{brownianLangevinNoise\_\allowbreak{}increment} \quad UnadjustedLangevin/\allowbreak{}BrownianSegmentsGlobal.lean, lines 24–30}
\begin{leancode}
\leanline{0}{theorem\ brownianLangevinNoise\_increment}
\leanline{4}{(W\ :\ ι\ →\ ℝ≥0\ →\ Ω\ →\ ℝ)\ (hW\ :\ ∀\ i,\ IsPreBrownianReal\ (W\ i)\ P)\ (a\ t\ :\ ℝ)\ :}
\leanline{4}{brownianLangevinNoise\ P\ W\ hW\ 0\ (a\ +\ t)\ -\ brownianLangevinNoise\ P\ W\ hW\ 0\ a\ =}
\leanline{6}{brownianLangevinNoise\ P\ W\ hW\ a\ t}
\end{leancode}
\leanlabel{How the Lean proof follows the argument}
\texttt{brownian\allowbreak{}Langevin\allowbreak{}Noise\_\allowbreak{}increment} is stated before the section variable \texttt{[Is\allowbreak{}Probability\allowbreak{}Measure P]} is introduced, so \(P\) is an arbitrary measure. The proof unfolds \texttt{brownian\allowbreak{}Langevin\allowbreak{}Noise}, simplifies \(0+s\) to \(s\) with \texttt{simp only [zero\_\allowbreak{}add]}, and closes the resulting identity between linear combinations of the three classes \(\beta(a+t)\), \(\beta(a)\), \(\beta(0)\) with the tactic \texttt{module}, which normalises both sides as linear combinations with real coefficients and compares them; this is the displayed computation.
\end{leanbox}

\section{\texorpdfstring{Shifting the Langevin integral equation}{Shifting the Langevin integral equation}}


\textbf{Claim.} Let \(P\) be a probability measure, \(\nabla U\) \(L\)-Lipschitz, \(X_0\in L^2\), \((W^i)\) pre-Brownian motions, \(T\in\mathbb R\), and let \(Y:\mathbb R\to L^2\) be continuous and satisfy \((\ast)\) on \([0,T]\). Let \(a\ge0\), \(h\ge0\) with \(a+h\le T\). Then for every \(t\in[0,h]\),


\[Y(a+t)-Y(a)=-\int_0^t\mathcal G\bigl(Y(a+s)\bigr)\,ds+N_a(t).\]


\textbf{Proof.} Fix \(t\in[0,h]\). Then \(a\in[0,T]\), because \(0\le a\le a+h\le T\), and \(a+t\in[0,T]\), because \(0\le a\le a+t\le a+h\le T\).




\emph{Increment formula.} The map \(s\mapsto\mathcal G(Y(s))\) is continuous, as the composition of the continuous \(Y\) with the Lipschitz map \(\mathcal G\); hence it is Bochner integrable on every bounded interval, and oriented integrals are additive: \(\int_0^{a+t}=\int_0^a+\int_a^{a+t}\). Writing \((\ast)\) at the times \(a+t\) and \(a\) and subtracting, the term \(X_0\) cancels and


\[Y(a+t)-Y(a)=-\int_a^{a+t}\mathcal G\bigl(Y(s)\bigr)\,ds+\bigl(N_0(a+t)-N_0(a)\bigr).\qquad(1)\]


This is the general increment identity for solutions of the integral equation \cite[Section~3.16.4]{LangevinLean2026ULASC}.




\emph{Change of variables.} Lebesgue measure is invariant under the translation \(r\mapsto a+r\), so


\[\int_a^{a+t}\mathcal G\bigl(Y(s)\bigr)\,ds=\int_0^t\mathcal G\bigl(Y(a+r)\bigr)\,dr.\qquad(2)\]


\emph{Conclusion.} Substitute (2) and the identity \(N_0(a+t)-N_0(a)=N_a(t)\) of the previous step into (1). \(\square\)




The conclusion says that \(r\mapsto Y(a+r)\) solves, on \([0,h]\), the segment equation with initial value \(Y(a)\) and input \(N_a\). The condition \(a+h\le T\) keeps both times \(a\) and \(a+t\) inside the interval on which \((\ast)\) is assumed.

\begin{leanbox}{Lean: \texttt{lp\_\allowbreak{}brownian\_\allowbreak{}integral\_\allowbreak{}shift}}
\leansource{\texttt{lp\_\allowbreak{}brownian\_\allowbreak{}integral\_\allowbreak{}shift} \quad UnadjustedLangevin/\allowbreak{}BrownianSegmentsGlobal.lean, lines 36–63}
\begin{leancode}
\leanline{0}{theorem\ lp\_brownian\_integral\_shift}
\leanline{4}{(U\ :\ EuclideanSpace\ ℝ\ ι\ →\ ℝ)\ (L\ :\ ℝ≥0)}
\leanline{4}{(hLip\ :\ LipschitzWith\ L\ (gradient\ U))}
\leanline{4}{(initial\ :\ Lp\ (EuclideanSpace\ ℝ\ ι)\ 2\ P)}
\leanline{4}{(W\ :\ ι\ →\ ℝ≥0\ →\ Ω\ →\ ℝ)\ (hW\ :\ ∀\ i,\ IsPreBrownianReal\ (W\ i)\ P)}
\leanline{4}{(Y\ :\ ℝ\ →\ Lp\ (EuclideanSpace\ ℝ\ ι)\ 2\ P)\ (hY\ :\ Continuous\ Y)}
\leanline{4}{(T\ :\ ℝ)}
\leanline{4}{(hYe\ :\ ∀\ t\ ∈\ Icc\ 0\ T,\ Y\ t\ -\ initial\ =}
\leanline{6}{-(∫\ s\ in\ 0..t,\ lpDrift\ P\ (gradient\ U)\ hLip\ (Y\ s))\ +}
\leanline{8}{brownianLangevinNoise\ P\ W\ hW\ 0\ t)}
\leanline{4}{(a\ h\ :\ ℝ)\ (ha\ :\ 0\ ≤\ a)\ (hh\ :\ 0\ ≤\ h)\ (hah\ :\ a\ +\ h\ ≤\ T)\ :}
\leanline{4}{∀\ t\ ∈\ Icc\ 0\ h,\ Y\ (a\ +\ t)\ -\ Y\ a\ =}
\leanline{6}{-(∫\ s\ in\ 0..t,\ lpDrift\ P\ (gradient\ U)\ hLip\ (Y\ (a\ +\ s)))\ +}
\leanline{8}{brownianLangevinNoise\ P\ W\ hW\ a\ t}
\end{leancode}
\leanlabel{How the Lean proof follows the argument}
In \texttt{lp\_\allowbreak{}brownian\_\allowbreak{}integral\_\allowbreak{}shift}, \texttt{h\allowbreak{}Y : Continuous Y} and \texttt{h\allowbreak{}Ye} are the continuity of \(Y\) and \((\ast)\); \texttt{ha}, \texttt{hh}, \texttt{hah} are \(a\ge0\), \(h\ge0\), \(a+h\le T\). The proof follows the three parts above.



\begin{enumerate}


\item \texttt{hai} and \texttt{hati} are the memberships \(a\in[0,T]\) and \(a+t\in[0,T]\), proved by \texttt{linarith}.

\item \texttt{hinc} is (1). It is the lemma \texttt{langevin\_\allowbreak{}integral\_\allowbreak{}increment} of the section on Euler limits, applied in the Banach space \(L^2\) to the Lipschitz map \texttt{lp\allowbreak{}Drift P (gradient U) h\allowbreak{}Lip} (its Lipschitz constant is \texttt{lp\allowbreak{}Drift\_\allowbreak{}lipschitz P (gradient U) h\allowbreak{}Lip}) and the input \texttt{brownian\allowbreak{}Langevin\allowbreak{}Noise P W h\allowbreak{}W 0}.

\item \texttt{hi} is (2). It comes from Mathlib's \texttt{interval\allowbreak{}Integral.\allowbreak{}integral\_\allowbreak{}comp\_\allowbreak{}add\_\allowbreak{}left}, which states \(\int_{a'}^{b'}f(c+x)\,dx=\int_{c+a'}^{c+b'}f(x)\,dx\), used with \(c=a\), \(a'=0\), \(b'=t\); \texttt{simpa only [add\_\allowbreak{}zero]} simplifies \(a+0\) to \(a\).

\item \texttt{rw [hi, brownian\allowbreak{}Langevin\allowbreak{}Noise\_\allowbreak{}increment] at hinc} substitutes (2) and the input identity into (1), after which \texttt{hinc} is the goal.

\end{enumerate}
\end{leanbox}

\section{\texorpdfstring{A restarted segment equals the global solution}{A restarted segment equals the global solution}}


\textbf{Claim.} Assume the hypotheses of the previous step on \(P,U,L,X_0,(W^i),Y,T\) and \((\ast)\). Let \(h\ge0\), \(n\in\mathbb N\) and \(X\in L^2\) with


\[X=Y(nh)\qquad\text{and}\qquad(n+1)h\le T.\]


Then \(S_n(X)(t)=Y(nh+t)\) for every \(t\in[0,h]\).




\textbf{Proof.} Put \(Z(t)=Y(nh+t)\). This curve is continuous, as the composition of \(Y\) with a translation. Apply the previous step with \(a=nh\ge0\) (since \(n\ge0\) and \(h\ge0\)); its condition \(a+h\le T\) is \(nh+h=(n+1)h\le T\). Using \(X=Y(nh)=Z(0)\), it gives


\[Z(t)-X=-\int_0^t\mathcal G\bigl(Z(s)\bigr)\,ds+N_{nh}(t),\qquad t\in[0,h].\]


By its specification, \(S_n(X)\) is continuous and satisfies the same equation with the same initial value \(X\) and the same input \(N_{nh}\) on \([0,h]\).




Two continuous solutions on \([0,h]\) of the integral equation with the same initial value and the same additive input coincide on \([0,h]\) \cite[Section~3.43.4]{LangevinLean2026ULASC}. We recall the argument. Let \(D(t)=S_n(X)(t)-Z(t)\). Subtracting the two equations, \(X\) and \(N_{nh}(t)\) cancel, so


\[D(t)=-\int_0^t\Bigl(\mathcal G\bigl(S_n(X)(s)\bigr)-\mathcal G\bigl(Z(s)\bigr)\Bigr)\,ds,\qquad t\in[0,h].\]


The integrand is continuous, so by the fundamental theorem of calculus \(D\) is differentiable on \([0,h)\) from the right, with derivative of norm \(\|\mathcal G(S_n(X)(t))-\mathcal G(Z(t))\|_2\le L\|D(t)\|_2\), and \(D(0)=0\). Grönwall's inequality gives \(\|D(t)\|_2\le\|D(0)\|_2e^{Lt}=0\) on \([0,h]\). Hence \(S_n(X)=Z\) on \([0,h]\). \(\square\)




Because the inputs cancel in \(D\), only continuity of the input is required, not differentiability. The statement holds for whichever solution was chosen in the definition of \(S_n(X)\).

\begin{leanbox}{Lean: \texttt{brownianSegmentFrom\_\allowbreak{}eq\_\allowbreak{}global}}
\leansource{\texttt{brownianSegmentFrom\_\allowbreak{}eq\_\allowbreak{}global} \quad UnadjustedLangevin/\allowbreak{}BrownianSegmentsGlobal.lean, lines 67–91}
\begin{leancode}
\leanline{0}{theorem\ brownianSegmentFrom\_eq\_global}
\leanline{4}{(U\ :\ EuclideanSpace\ ℝ\ ι\ →\ ℝ)\ (L\ :\ ℝ≥0)}
\leanline{4}{(hLip\ :\ LipschitzWith\ L\ (gradient\ U))}
\leanline{4}{(initial\ initialAt\ :\ Lp\ (EuclideanSpace\ ℝ\ ι)\ 2\ P)}
\leanline{4}{(W\ :\ ι\ →\ ℝ≥0\ →\ Ω\ →\ ℝ)\ (hW\ :\ ∀\ i,\ IsPreBrownianReal\ (W\ i)\ P)}
\leanline{4}{(Y\ :\ ℝ\ →\ Lp\ (EuclideanSpace\ ℝ\ ι)\ 2\ P)\ (hY\ :\ Continuous\ Y)}
\leanline{4}{(T\ :\ ℝ)}
\leanline{4}{(hYe\ :\ ∀\ t\ ∈\ Icc\ 0\ T,\ Y\ t\ -\ initial\ =}
\leanline{6}{-(∫\ s\ in\ 0..t,\ lpDrift\ P\ (gradient\ U)\ hLip\ (Y\ s))\ +}
\leanline{8}{brownianLangevinNoise\ P\ W\ hW\ 0\ t)}
\leanline{4}{(h\ :\ ℝ)\ (hh\ :\ 0\ ≤\ h)\ (n\ :\ ℕ)}
\leanline{4}{(hstart\ :\ initialAt\ =\ Y\ ((n\ :\ ℝ)\ *\ h))}
\leanline{4}{(hend\ :\ ((n\ :\ ℝ)\ +\ 1)\ *\ h\ ≤\ T)\ :}
\leanline{4}{EqOn\ (brownianSegmentFrom\ P\ U\ L\ hLip\ W\ hW\ h\ hh\ n\ initialAt)}
\leanline{6}{(fun\ t\ ↦\ Y\ ((n\ :\ ℝ)\ *\ h\ +\ t))\ (Icc\ 0\ h)}
\end{leancode}
\leanlabel{How the Lean proof follows the argument}
\texttt{brownian\allowbreak{}Segment\allowbreak{}From\_\allowbreak{}eq\_\allowbreak{}global} has the arguments \texttt{initial} (\(X_0\)) and \texttt{initial\allowbreak{}At} (\(X\)), the hypotheses \texttt{hstart : initial\allowbreak{}At = Y ((n : \ensuremath{\mathbb R}) * h)} and \texttt{hend : ((n : \ensuremath{\mathbb R}) + 1) * h \ensuremath{\le} T}, and the conclusion \texttt{Eq\allowbreak{}On (brownian\allowbreak{}Segment\allowbreak{}From P U L h\allowbreak{}Lip W h\allowbreak{}W h hh n initial\allowbreak{}At) (fun t \ensuremath{\mapsto} Y ((n : \ensuremath{\mathbb R}) * h + t)) (Icc 0 h)}; \texttt{Eq\allowbreak{}On f g s} means \(f(t)=g(t)\) for all \(t\in s\).



\begin{enumerate}


\item \texttt{hs} is the specification \texttt{brownian\allowbreak{}Segment\allowbreak{}From\_\allowbreak{}spec} of \(S_n(X)\).

\item \texttt{he} is the shifted equation of the previous step, \texttt{lp\_\allowbreak{}brownian\_\allowbreak{}integral\_\allowbreak{}shift}, at \(a=nh\); the side condition \(nh+h\le T\) is derived from \texttt{hend} by \texttt{nlinarith}.

\item \texttt{apply lp\_\allowbreak{}langevin\_\allowbreak{}integral\_\allowbreak{}unique …} invokes the uniqueness theorem of the section on uniqueness with the two continuous curves (continuity of \(S_n(X)\) is \texttt{hs.\allowbreak{}1}; continuity of \(Z\) is \texttt{h\allowbreak{}Y.\allowbreak{}comp (continuous\_\allowbreak{}const.\allowbreak{}add continuous\_\allowbreak{}id)}) and the equation \texttt{hs.\allowbreak{}2.\allowbreak{}2.\allowbreak{}1} of \(S_n(X)\). The remaining goal is the equation of \(Z\) with initial value \texttt{initial\allowbreak{}At}; \texttt{rw [hstart]} replaces it by \(Y(nh)\) and \texttt{exact he t ht} concludes.

\end{enumerate}
\end{leanbox}

\section{\texorpdfstring{The knots lie on the global solution}{The knots lie on the global solution}}


\textbf{Claim.} Assume the hypotheses of the second step on \(P,U,L,X_0,(W^i),Y,T\) and \((\ast)\), and in addition \(T\ge0\). Let \(h\ge0\) and \(n\in\mathbb N\) with \(nh\le T\). Then \(K_n=Y(nh)\).




\textbf{Proof.} \emph{Initial value.} Since \(T\ge0\), the time \(0\) lies in \([0,T]\), and \((\ast)\) at \(t=0\) reads \(Y(0)-X_0=-\int_0^0\mathcal G(Y(s))\,ds+N_0(0)=0+0\), using that an integral over a degenerate interval vanishes and \(N_0(0)=0\). Hence \(Y(0)=X_0\).




We prove by induction on \(n\) the statement: if \(nh\le T\) then \(K_n=Y(nh)\).




\emph{Base case \(n=0\).} \(K_0=X_0=Y(0)=Y(0\cdot h)\).




\emph{Inductive step.} Assume the statement for \(n\) and let \((n+1)h\le T\). Since \(h\ge0\), \(nh\le nh+h=(n+1)h\le T\), so the induction hypothesis gives \(K_n=Y(nh)\). The previous step, applied with \(X=K_n\), gives \(S_n(K_n)(t)=Y(nh+t)\) for \(t\in[0,h]\); at \(t=h\in[0,h]\),


\[K_{n+1}=S_n(K_n)(h)=Y(nh+h)=Y\bigl((n+1)h\bigr).\]


This completes the induction. \(\square\)

\begin{leanbox}{Lean: \texttt{brownianKnots\_\allowbreak{}eq\_\allowbreak{}global}}
\leansource{\texttt{brownianKnots\_\allowbreak{}eq\_\allowbreak{}global} \quad UnadjustedLangevin/\allowbreak{}BrownianSegmentsGlobal.lean, lines 95–124}
\begin{leancode}
\leanline{0}{theorem\ brownianKnots\_eq\_global}
\leanline{4}{(U\ :\ EuclideanSpace\ ℝ\ ι\ →\ ℝ)\ (L\ :\ ℝ≥0)}
\leanline{4}{(hLip\ :\ LipschitzWith\ L\ (gradient\ U))}
\leanline{4}{(initial\ :\ Lp\ (EuclideanSpace\ ℝ\ ι)\ 2\ P)}
\leanline{4}{(W\ :\ ι\ →\ ℝ≥0\ →\ Ω\ →\ ℝ)\ (hW\ :\ ∀\ i,\ IsPreBrownianReal\ (W\ i)\ P)}
\leanline{4}{(Y\ :\ ℝ\ →\ Lp\ (EuclideanSpace\ ℝ\ ι)\ 2\ P)\ (hY\ :\ Continuous\ Y)}
\leanline{4}{(T\ :\ ℝ)\ (hT\ :\ 0\ ≤\ T)}
\leanline{4}{(hYe\ :\ ∀\ t\ ∈\ Icc\ 0\ T,\ Y\ t\ -\ initial\ =}
\leanline{6}{-(∫\ s\ in\ 0..t,\ lpDrift\ P\ (gradient\ U)\ hLip\ (Y\ s))\ +}
\leanline{8}{brownianLangevinNoise\ P\ W\ hW\ 0\ t)}
\leanline{4}{(h\ :\ ℝ)\ (hh\ :\ 0\ ≤\ h)\ (n\ :\ ℕ)\ (hn\ :\ (n\ :\ ℝ)\ *\ h\ ≤\ T)\ :}
\leanline{4}{brownianKnots\ P\ U\ L\ hLip\ W\ hW\ h\ hh\ initial\ n\ =\ Y\ ((n\ :\ ℝ)\ *\ h)}
\end{leancode}
\leanlabel{How the Lean proof follows the argument}
In \texttt{brownian\allowbreak{}Knots\_\allowbreak{}eq\_\allowbreak{}global} the additional hypotheses are \texttt{h\allowbreak{}T : 0 \ensuremath{\le} T} and \texttt{hn : (n : \ensuremath{\mathbb R}) * h \ensuremath{\le} T}.



\begin{enumerate}


\item \texttt{h\allowbreak{}Y0 : Y 0 = initial} evaluates \texttt{h\allowbreak{}Ye} at \(0\in[0,T]\) and simplifies with \texttt{interval\allowbreak{}Integral.\allowbreak{}integral\_\allowbreak{}same} (the integral over \([0,0]\) is \(0\)), \texttt{brownian\allowbreak{}Langevin\allowbreak{}Noise\_\allowbreak{}zero} and \texttt{sub\_\allowbreak{}eq\_\allowbreak{}zero}.

\item \texttt{induction n with} performs the induction; since \texttt{hn} mentions \(n\), it is generalised, so the induction hypothesis \texttt{ih} has the form \(nh\le T\Rightarrow K_n=Y(nh)\). The \texttt{zero} case unfolds \texttt{brownian\allowbreak{}Knots} and uses \texttt{h\allowbreak{}Y0}.

\item In the \texttt{succ} case, \texttt{hprev} derives \(nh\le T\) from \texttt{hn} (\texttt{push\_\allowbreak{}cast} moves the coercion of \(n+1\) into the reals and \texttt{nlinarith} uses \(h\ge0\)); \texttt{hstart := ih hprev} is \(K_n=Y(nh)\); \texttt{he} is the previous step with \texttt{initial\allowbreak{}At} equal to the knot; \texttt{heh} evaluates it at \(t=h\). After \texttt{change} exhibits \(K_{n+1}\) as \(S_n(K_n)(h)\), \texttt{simpa} rewrites \(nh+h\) as \((n+1)h\) and closes the goal.

\end{enumerate}
\end{leanbox}

\section{\texorpdfstring{Every reference segment lies on the global solution}{Every reference segment lies on the global solution}}


\textbf{Claim.} Assume the hypotheses of the second step on \(P,U,L,X_0,(W^i),Y,T\) and \((\ast)\), and \(T\ge0\). Let \(h\ge0\) and \(n\in\mathbb N\) with \((n+1)h\le T\). Then


\[Y^{(n)}(t)=Y(nh+t)\qquad\text{for every }t\in[0,h].\]


\textbf{Proof.} Since \(h\ge0\), \(nh\le(n+1)h\le T\), so the previous step gives \(K_n=Y(nh)\). The third step, applied with \(X=K_n\) and the hypothesis \((n+1)h\le T\), gives \(S_n(K_n)(t)=Y(nh+t)\) for \(t\in[0,h]\), and \(Y^{(n)}=S_n(K_n)\) by definition. \(\square\)




Consequently, if \(Y(s)\) has a given law \(\pi\) for every \(s\in[0,T]\), then so does every \(Y^{(n)}(t)\) with \(t\in[0,h]\) and \((n+1)h\le T\). This is how the stationarity hypothesis of the reference construction is reduced to a statement about a single global solution.

\begin{leanbox}{Lean: \texttt{brownianSegments\_\allowbreak{}eq\_\allowbreak{}global}}
\leansource{\texttt{brownianSegments\_\allowbreak{}eq\_\allowbreak{}global} \quad UnadjustedLangevin/\allowbreak{}BrownianSegmentsGlobal.lean, lines 128–144}
\begin{leancode}
\leanline{0}{theorem\ brownianSegments\_eq\_global}
\leanline{4}{(U\ :\ EuclideanSpace\ ℝ\ ι\ →\ ℝ)\ (L\ :\ ℝ≥0)}
\leanline{4}{(hLip\ :\ LipschitzWith\ L\ (gradient\ U))}
\leanline{4}{(initial\ :\ Lp\ (EuclideanSpace\ ℝ\ ι)\ 2\ P)}
\leanline{4}{(W\ :\ ι\ →\ ℝ≥0\ →\ Ω\ →\ ℝ)\ (hW\ :\ ∀\ i,\ IsPreBrownianReal\ (W\ i)\ P)}
\leanline{4}{(Y\ :\ ℝ\ →\ Lp\ (EuclideanSpace\ ℝ\ ι)\ 2\ P)\ (hY\ :\ Continuous\ Y)}
\leanline{4}{(T\ :\ ℝ)\ (hT\ :\ 0\ ≤\ T)}
\leanline{4}{(hYe\ :\ ∀\ t\ ∈\ Icc\ 0\ T,\ Y\ t\ -\ initial\ =}
\leanline{6}{-(∫\ s\ in\ 0..t,\ lpDrift\ P\ (gradient\ U)\ hLip\ (Y\ s))\ +}
\leanline{8}{brownianLangevinNoise\ P\ W\ hW\ 0\ t)}
\leanline{4}{(h\ :\ ℝ)\ (hh\ :\ 0\ ≤\ h)\ (n\ :\ ℕ)\ (hend\ :\ ((n\ :\ ℝ)\ +\ 1)\ *\ h\ ≤\ T)\ :}
\leanline{4}{EqOn\ (brownianSegments\ P\ U\ L\ hLip\ W\ hW\ h\ hh\ initial\ n)}
\leanline{6}{(fun\ t\ ↦\ Y\ ((n\ :\ ℝ)\ *\ h\ +\ t))\ (Icc\ 0\ h)}
\end{leancode}
\leanlabel{How the Lean proof follows the argument}
\texttt{brownian\allowbreak{}Segments\_\allowbreak{}eq\_\allowbreak{}global} has the hypotheses \texttt{h\allowbreak{}T : 0 \ensuremath{\le} T} and \texttt{hend : ((n : \ensuremath{\mathbb R}) + 1) * h \ensuremath{\le} T}. The proof first obtains \texttt{hkn}, the knot identity \(K_n=Y(nh)\), from \texttt{brownian\allowbreak{}Knots\_\allowbreak{}eq\_\allowbreak{}global}, whose side condition \(nh\le T\) is proved by \texttt{nlinarith}; it then applies \texttt{brownian\allowbreak{}Segment\allowbreak{}From\_\allowbreak{}eq\_\allowbreak{}global} with \texttt{initial\allowbreak{}At} equal to the knot, the identity \texttt{hkn} and \texttt{hend}. Since \texttt{brownian\allowbreak{}Segments} unfolds to \texttt{brownian\allowbreak{}Segment\allowbreak{}From} at the knot, this is the conclusion.
\end{leanbox}

\begin{thebibliography}{ULA-LSI}
\bibitem[ULA-SC]{LangevinLean2026ULASC}
LangevinLean Project. \emph{Unadjusted Langevin under Strong Convexity}. LL-01, 2026. Draft report, content revision d34ead06987c. \url{https://odakyildiz.github.io/LangevinLean/docs/pdf/UnadjustedLangevin/FullProof.pdf}.
\end{thebibliography}

\end{document}
